\documentclass[a4paper, serif, 10pt]{moderncv}

%% moderncv
% color and style
\moderncvcolor{black}
\moderncvstyle{banking}

% renew moderncv command to adapt for CJK colon
\renewcommand*{\cvitem}[3][.25em]{%
  \ifstrempty{#2}{}{\hintstyle{#2}：}{#3}%
  \par\addvspace{#1}}

\renewcommand*{\cvitemwithcomment}[4][.25em]{%
  \savebox{\cvitemwithcommentmainbox}{\ifstrempty{#2}{}{\hintstyle{#2}：}#3}%
  \setlength{\cvitemwithcommentmainlength}{\widthof{\usebox{\cvitemwithcommentmainbox}}}%
  \setlength{\cvitemwithcommentcommentlength}{\maincolumnwidth-\separatorcolumnwidth-\cvitemwithcommentmainlength}%
  \begin{minipage}[t]{\cvitemwithcommentmainlength}\usebox{\cvitemwithcommentmainbox}\end{minipage}%
  \hfill% fill of \separatorcolumnwidth
  \begin{minipage}[t]{\cvitemwithcommentcommentlength}\raggedleft\small\itshape#4\end{minipage}%
  \par\addvspace{#1}}

%% page layout/margins
\usepackage[top=2.5cm, bottom=2.5cm, left=1.5cm, right=1.5cm]{geometry}

\nopagenumbers{}

%% Babel config for English language
\usepackage[english]{babel}

%% fontspec
\usepackage{fontspec}

\IfFontExistsTF{Linux Libertine}{
  \setmainfont[Ligatures={TeX, Common}, Numbers=Lining, ItalicFont=Linux Libertine]{Linux Libertine}
}{}
\IfFontExistsTF{Linux Libertine O}{
  \setmainfont[Ligatures={TeX, Common}, Numbers=Lining, ItalicFont=Linux Libertine O]{Linux Libertine O}
}{}

%% CTeX
% CJK support, used to show CJK characters in the resume
%
% - heading=false: disable ctex heading
% - punct=kaiming: use kaiming punctuations styles for CJK
% - scheme=plain: use plain scheme, do not override \normalsize font size
% - space=auto: space settings for CJK characters
%
% CJK fonts are resolved per locale language, e.g. Noto Serif/Sans CJK JP
% for Japanese resumes, with fallback to the next region if not installed
%
% ref:
% - http://ctan.mirrorcatalogs.com/language/chinese/ctex/ctex.pdf
\usepackage[UTF8, heading=false, punct=kaiming, scheme=plain, space=auto]{ctex}

\IfFontExistsTF{Noto Serif CJK TC}{
  \setCJKmainfont{Noto Serif CJK TC}
}{}
\IfFontExistsTF{Noto Sans CJK TC}{
  \setCJKsansfont{Noto Sans CJK TC}
}{}

%% line spacing
\usepackage{setspace}
\setstretch{1.125}

%% URL styling - use normal text instead of monospace
\urlstyle{same}

%% Patch \httplink and \httpslink to handle full URLs
% The original moderncv macros always prepend http:// or https:// to URLs.
% This causes issues when the URL already has a protocol (e.g., https://example.com)
% resulting in malformed URLs like https://https://example.com.
% This patch checks if the URL already contains "://" and uses it directly if so.
% Note: We use \str_set:Nx (x-expansion) to expand macros like \@homepage first.
\ExplSyntaxOn
\str_new:N \g_yamlresume_url_str

\RenewDocumentCommand{\httpslink}{O{}m}{
  \str_set:Nx \g_yamlresume_url_str {#2}
  \str_if_in:NnTF \g_yamlresume_url_str {://}
    { \url{#2} }
    { \url{https://#2} }
}

\RenewDocumentCommand{\httplink}{O{}m}{
  \str_set:Nx \g_yamlresume_url_str {#2}
  \str_if_in:NnTF \g_yamlresume_url_str {://}
    { \url{#2} }
    { \url{http://#2} }
}
\ExplSyntaxOff

\name{林彥廷}{}
\title{資深資安工程師 | 紅隊演練與事件回應}
\phone[mobile]{+886 2 2712 3456}
\email{yenting.lin@cybersec.tw}
\homepage{https://www.yen-ting-lin.dev}

\address{中國臺灣，台北市，台北市，台北市大安區復興南路一段 123 號 8 樓，106}{}{}

\extrainfo{{\small \faGithub}\ \href{https://github.com/yentinglin-cyber}{@yentinglin-cyber} {} {} {} • {} {} {} 
{\small \faLinkedin}\ \href{https://www.linkedin.com/in/yentinglin-cyber}{@yentinglin-cyber} {} {} {} • {} {} {} 
{\small \faStackOverflow}\ \href{https://stackoverflow.com/users/12345678/yenting-lin}{@yenting-lin}}

\begin{document}

\maketitle

\section{簡介}

\cvline{}{\begin{itemize}
\item 具備超過 8 年資訊安全經驗，專注於滲透測試、紅隊演練、事件回應與資安維運。
\item 熟悉 MITRE ATT\&CK、OWASP Top 10、NIST CSF 與 ISO 27001 等框架，能將攻擊手法轉化為可執行的防禦策略。
\item 擅長使用 Python、Go 與 Shell 開發自動化工具，提升弱點管理與威脅偵測效率。
\item 具備跨部門溝通與專案管理能力，曾帶領 4 人小組完成金融業零信任架構導入。
\end{itemize}}
\section{教育背景}

\cventry{2014年9月–2016年6月}
        {碩士，資訊工程學系，成績：4.1/4.3}
        {國立臺灣大學}
        {\url{https://www.ntu.edu.tw/}}
        {}
        {\begin{itemize}
\item 研究主題聚焦於入侵偵測與惡意流量分析
\item 參與實驗室與企業合作之資安專案，建立威脅情報分析流程
\item 擔任課程助教，協助教授計算機安全與密碼學課程
\end{itemize}
\textbf{課程}：計算機安全、密碼學、網路安全、作業系統、資料庫系統、分散式系統}

\cventry{2010年9月–2014年6月}
        {學士，資訊工程學系，成績：3.9/4.3}
        {國立交通大學}
        {\url{https://www.nctu.edu.tw/}}
        {}
        {\begin{itemize}
\item 建立程式設計與計算機網路基礎
\item 參加校內資安社團，開始接觸 CTF 競賽
\end{itemize}
\textbf{課程}：資料結構、演算法、計算機組織、作業系統、網路程式設計}
\section{工作經歷}

\cventry{2021年3月至今}
        {資深資安工程師}
        {藍盾資訊安全股份有限公司}
        {\url{https://www.blueshield-sec.tw/}}
        {}
        {\begin{itemize}
\item 主導滲透測試與紅隊演練專案，涵蓋金融、電商與政府單位，年均完成 20 餘件專案
\item 建立以 MITRE ATT\&CK 為基礎的攻擊模擬與偵測規則，提升 SIEM 告警準確率 35\%
\item 規劃與執行事件回應流程，協調跨部門進行入侵調查、橫向移動分析與復原
\item 帶領 4 人資安小組，負責專案排程、技術指導與客戶溝通
\item 推動 ISO 27001 與 NIST CSF 控制措施落地，協助客戶通過第三方稽核
\end{itemize}
\textbf{關鍵字}：滲透測試、紅隊演練、事件回應、MITRE ATT\&CK、SIEM、弱點管理、ISO 27001}

\cventry{2018年4月–2021年2月}
        {資安工程師}
        {雲峰科技股份有限公司}
        {\url{https://www.cloudpeak-tech.tw/}}
        {}
        {\begin{itemize}
\item 負責企業客戶資安健診、弱點掃描與滲透測試，提供修補建議與風險排序
\item 建置與維運 ELK 日誌分析平台，設計威脅偵測規則與儀表板
\item 協助客戶處理資安事件，進行惡意程式初步分析與數位鑑識
\item 撰寫資安報告與 SOP，並對客戶進行資安意識教育訓練
\end{itemize}
\textbf{關鍵字}：弱點掃描、滲透測試、ELK、數位鑑識、資安健診、教育訓練}

\cventry{2016年7月–2018年3月}
        {網路安全分析師}
        {台灣數位金融股份有限公司}
        {\url{https://www.twdigitalfinance.tw/}}
        {}
        {\begin{itemize}
\item 監控網路與端點安全事件，進行告警分流、調查與升級處理
\item 維運防火牆、WAF 與 IDS/IPS，定期檢視規則與調整政策
\item 參與行動銀行 App 安全檢測，協助修補 OWASP Top 10 相關弱點
\item 建立資安事件通報流程與每月營運報表
\end{itemize}
\textbf{關鍵字}：事件監控、WAF、IDS/IPS、OWASP Top 10、行動安全、通報流程}
\section{語言}

\cvline{中文}{母語或雙語水平 \hfill \textbf{關鍵字}：母語}
\cvline{英語}{完全專業水平 \hfill \textbf{關鍵字}：TOEIC 950、IELTS 7.5}
\cvline{日語}{初級水平 \hfill \textbf{關鍵字}：JLPT N4}
\section{專業技能}

\cvline{滲透測試與紅隊}{專家 \hfill \textbf{關鍵字}：OWASP Top 10、Burp Suite、Metasploit、Nmap、Kali Linux、社交工程}
\cvline{事件回應與鑑識}{高級 \hfill \textbf{關鍵字}：SIEM、Splunk、ELK、數位鑑識、惡意程式分析、MITRE ATT\&CK}
\cvline{程式與自動化}{高級 \hfill \textbf{關鍵字}：Python、Go、Bash、PowerShell、SQL、REST API}
\cvline{雲端與容器安全}{中級 \hfill \textbf{關鍵字}：AWS、Azure、GCP、IAM、Kubernetes、容器安全}
\cvline{資安治理與合規}{高級 \hfill \textbf{關鍵字}：ISO 27001、NIST CSF、風險管理、合規稽核、資安政策}
\section{榮譽}

\cventry{2023年12月}
        {藍盾資訊安全股份有限公司}
        {年度最佳資安專案獎}
        {}
        {}
        {\begin{itemize}
\item 表揚於金融業紅隊演練專案中發現重大風險並提出有效改善方案
\item 專案協助客戶降低橫向移動風險並強化偵測能力
\end{itemize}}

\cventry{2019年8月}
        {台灣資安社群聯盟}
        {台灣資安 CTF 競賽第一名}
        {}
        {}
        {\begin{itemize}
\item 帶領 4 人團隊於 Web、Pwn 與 Reverse 題型中取得佳績
\item 負責解題策略規劃與弱點利用腳本開發
\end{itemize}}
\section{證書}

\cventry{2022年6月}
        {ISC2}
        {CISSP 國際資訊系統安全認證}
        {\url{https://www.isc2.org/certifications/cissp}}
        {}
        {}

\cventry{2020年11月}
        {Offensive Security}
        {OSCP 滲透測試認證}
        {\url{https://www.offsec.com/courses/pen-200/}}
        {}
        {}

\cventry{2021年9月}
        {Amazon Web Services}
        {AWS Certified Security - Specialty}
        {\url{https://aws.amazon.com/certification/certified-security-specialty/}}
        {}
        {}

\cventry{2019年5月}
        {EC-Council}
        {CEH 道德駭客認證}
        {\url{https://www.eccouncil.org/programs/certified-ethical-hacker-ceh/}}
        {}
        {}
\section{出版刊物}

\cventry{2022年11月}
        {以 MITRE ATT\&CK 為基礎的供應鏈攻擊偵測方法}
        {台灣資訊安全期刊}
        {\url{https://www.taiwan-infosec-journal.tw/articles/attck-supply-chain}}
        {}
        {\begin{itemize}
\item 提出結合 ATT\&CK 技術與 SIEM 規則的偵測框架
\item 以實際事件案例驗證偵測覆蓋率與誤報率改善
\end{itemize}}
\section{項目}

\cventry{2022年1月–2023年6月}
        {以 Python 與 Go 開發的紅隊演練輔助平台，整合偵察、弱點驗證與報告產出。}
        {紅隊演練自動化平台}
        {\url{https://github.com/yentinglin-cyber/redteam-automation}}
        {}
        {\begin{itemize}
\item 整合 Nmap、Burp Suite 與自製腳本，自動化執行偵察與弱點驗證
\item 建立 ATT\&CK 技術對應標籤，協助團隊快速彙整攻擊路徑與防禦建議
\item 將報告產出時間由 2 天縮短至 4 小時
\end{itemize}
\textbf{關鍵字}：自動化、紅隊演練、Python、Go、ATT\&CK}

\cventry{2021年8月–2022年12月}
        {協助金融客戶規劃與導入零信任架構，涵蓋身分驗證、裝置信任與微隔離。}
        {金融業零信任架構導入}
        {}
        {}
        {\begin{itemize}
\item 盤點關鍵資產與高風險存取路徑，制定分階段導入計畫
\item 導入多因子驗證、條件式存取與網路微隔離控制
\item 建立監控指標與異常行為告警，提升橫向移動偵測能力
\end{itemize}
\textbf{關鍵字}：零信任、身分驗證、微隔離、條件式存取、金融安全}
\section{興趣愛好}

\cvline{CTF 競賽}{Web、Pwn、Reverse、Crypto}
\cvline{戶外活動}{登山、馬拉松、攝影}
\section{志願服務}

\cventry{2019年1月–2023年12月}
        {資安志工講師}
        {台灣資安教育推廣協會}
        {\url{https://www.taiwan-cyber-edu.org/}}
        {}
        {\begin{itemize}
\item 至大專院校與高中分享資安職涯與基礎防護觀念
\item 協助辦理資安工作坊與 CTF 新手教學
\item 參與偏鄉資訊教育營隊，教授網路安全與密碼學入門
\end{itemize}}
    
\end{document}